\documentclass{article}
\usepackage{amsmath, amssymb, amsfonts}
\usepackage[utf8]{inputenc}
\usepackage[spanish]{babel}
\usepackage{hyperref}

\title{Taller II}
\author{Iván Palomo}
\date{20 de febrero de 2026}

\begin{document}
	\maketitle
	\tableofcontents
	\newpage
	
	\section{$|x^2-1|<x+1$}
	
	
	
	\medskip
	Como $|x^2-1|\ge 0$, para que se cumpla
	$|x^2-1|<x+1$ es necesario que $x+1>0$, es decir,
	
	$$
	x>-1.
	$$
	
	Si $x+1>0$, entonces
	
	
	$$|x^2-1|<x+1 \iff -(x+1)<x^2-1<x+1.$$
	
	
	Esto equivale a resolver:
	
	$$
	\begin{cases}
		x^2-1<x+1,\\
		-(x+1)<x^2-1.
	\end{cases}
	$$
	
	Resolvemos cada desigualdad por separado:
	\begin{align*}
		x^2-1<x+1
		&\iff x^2-x-2<0\\
		&\iff (x-2)(x+1)<0\\
		&\iff -1<x<2.
	\end{align*}
	
	\begin{align*}
		-(x+1)<x^2-1
		&\iff -x-1<x^2-1\\
		&\iff 0<x^2+x\\
		&\iff x(x+1)>0\\
		&\iff x<-1 \ \text{o}\ x>0.
	\end{align*}
	
	Juntando $-1<x<2$, con $(x<-1)\lor(x>0)$ y además $x>-1$, obtenemos
	
	
	$$
	\boxed{ x\in(0,2)}
	$$
	
	
	\section{$x^3-x^2+2x-2<0$}
	
	Factorizando:
	
	$$
	x^3-x^2+2x-2 =(x-1)(x^2+2)
	$$
	
	Observamos que:
	
	$$
	x^2+2>0 \quad \forall x\in\mathbb{R}
	$$
	
	Por tanto, el signo depende solo de $(x-1)$.
	
	$$
	(x-1)(x^2+2)<0 \iff x-1<0
	$$
	
	$$
	\boxed{x\in(-\infty , 1)}
	$$
	
	
	\section{$|x+3|+|x-1|\ge 5$}
	
	Los puntos críticos son $x=-3$ y $x=1$.
	
	\subsection*{Caso 1: $x\le -3$}
	
	$$
	|x+3|=-(x+3),\quad |x-1|=-(x-1)
	$$
	
	$$
	|x+3|+|x-1|=-2x-2
	$$
	
	$$
	-2x-2\ge5 \iff x\le -\frac{7}{2}
	$$
	
	\subsection*{Caso 2: $-3\le x\le 1$}
	
	$$
	|x+3|+|x-1|=4
	$$
	
	No hay soluciones en este intervalo, es decir,
	
	$$
	\varnothing
	$$
	
	\subsection*{Caso 3: $x\ge 1$}
	
	$$
	|x+3|=x+3,\quad |x-1|=x-1
	$$
	
	$$
	|x+3|+|x-1|=2x+2
	$$
	
	$$
	2x+2\ge5 \iff x\ge \frac{3}{2}
	$$
	
	
	
	Tomando la unión de los casos válidos:
	
	$$
	\boxed{x\in(-\infty,-\tfrac{7}{2}] \cup [\tfrac{3}{2},\infty)}
	$$
	
	
	
\end{document}